% ---------------------------------------------------------------------------
% Kapitel 3: Grundlagen und verwandte Arbeiten
% Kurzfassung vom 04.10.2026, neu geschrieben aus der Langfassung
% (thesis/langfassung/kapitel03.tex, Confluence-Seite 34373645).
% Tabelle 3.1 fuehrt die frueheren Tabellen 3.1 und 3.2 zusammen.
% ---------------------------------------------------------------------------

\chapter{Grundlagen und verwandte Arbeiten}
\label{kap:grundlagen}

Drei Stränge der Literatur tragen diese Arbeit. Die E-Portfolio-Literatur sagt, was ein Portfolio ist und woraus es besteht, die Akzeptanzforschung, warum ein freiwillig gefülltes System an seiner Bedienbarkeit scheitern kann, und die Usability-Literatur, woran sich Bedienbarkeit bemisst. Diese Arbeitsteilung bestimmt, welche Anforderung sich auf einen Beleg stützen kann und welche eine eigene Ableitung bleibt. Das Kapitel bearbeitet die erste Teilfrage und liefert in Abschnitt~\ref{sec:usability} die Kriterien der zweiten.

\section{E-Portfolio: Begriffe und Kategorien}
\label{sec:eportfolio}

Lorenzo und Ittelson (2005) bestimmen ein E-Portfolio als digitalisierte Sammlung von Artefakten, die eine Einzelperson, eine Gruppe oder eine Institution repräsentiert. Es könne zugleich als Verwaltungswerkzeug dienen, das die Arbeiten ordnet und steuert, wer sie sehen darf. Sichtbarkeitssteuerung ist damit im Begriff angelegt und keine Zutat des Baus. Die Autoren unterscheiden studentische, lehrbezogene und institutionelle Portfolios und halten fest, dass sich die ersten beiden auf der breitesten Ebene zu einem institutionellen zusammenführen lassen. Die vorliegende Arbeit liest die dritte Kategorie deshalb als Aggregationsebene: Ein Portal, das studentische Arbeiten sammelt und zeigt, ist Ort ihrer Entstehung und zugleich Aggregat, das die Einrichtung vertritt.

Für das institutionelle Portfolio geben Lorenzo und Ittelson eine Bestimmung wieder, nach der es eine fokussierte Auswahl authentischer Arbeiten ist, die Rechenschaft der Einrichtung belegt (Kahn, 2001, zitiert nach Lorenzo \& Ittelson, 2005). Daraus folgen zwei Festlegungen. Eine fokussierte Auswahl verlangt einen Auswahlschritt, also eine redaktionelle Freigabe (Abschnitt~\ref{sec:workflow}). Und weil die Autoren die Leserschaft veröffentlichter Unterlagen einer Einrichtung als offen beschreiben, legt die vorliegende Arbeit den Zugang ohne Anmeldung fest; die Übertragung auf studentische Arbeiten ist ihr eigener Schritt. Ebenso eine Setzung ist die Einordnung als Präsentationsportfolio, das fertige Ergebnisse zeigt und keine prozessbegleitenden Artefakte.

Was ein hochschulisches E-Portfolio enthalten sollte, beantwortet Mahasneh (2020). Aus einer Literaturdurchsicht und Interviews mit 20 Fachleuten entsteht ein Modell aus zehn Komponenten: Biografie, Kursplan, Berichte und Forschungsarbeiten, Hausaufgaben, Projekte und Experimente, Aktivitäten, Zusammenfassungen, wissenschaftliches Material, Audio- und Videoclips sowie Leistungsproben. 90 Studierende bewerteten es mit einem Mittelwert von 4,58 auf einer fünfstufigen Skala; die Aussage, ein E-Portfolio helfe beim Veröffentlichen eigener Erfahrungen, erreicht 4,96 bei geringer Streuung. Die vorliegende Arbeit ordnet die Komponenten nach ihrem Bezugspunkt, was die Quelle nicht tut: Personenbezogen ist allein die Biografie, an eine Lehrveranstaltung gebunden sind Kursplan, Hausaufgaben und Aktivitäten, die übrigen sechs sind artefaktbezogen. Daraus folgt die Trennung von Profilseite und Projekt, und die drei veranstaltungsgebundenen Komponenten entfallen mit der Abgrenzung gegenüber dem Lernmanagementsystem.

Beide Quellen haben Grenzen. Lorenzo und Ittelson (2005) legen einen Übersichtsbericht mit Momentaufnahmen von 2005 vor, Mahasneh (2020) erhebt ohne Kontrollgruppe an einem einzigen College in Jordanien. Das Modell trägt deshalb als Inhaltsraster, nicht als Wirkungsnachweis. Aus einem Whitepaper übernehmen Lorenzo und Ittelson zudem die Frage, ob ein System benutzerfreundlich genug für eine leichte Übernahme sei (ePortConsortium, 2003, zitiert nach Lorenzo \& Ittelson, 2005); die E-Portfolio-Literatur stellt die Akzeptanzfrage also, beantwortet sie aber nicht.

\section{Technologieakzeptanzmodell (TAM)}
\label{sec:tam}

Davis (1989) erklärt die Annahme eines Systems mit zwei Größen. Die wahrgenommene Nützlichkeit (PU) ist der Grad, zu dem eine Person glaubt, dass das System ihre Arbeitsleistung steigert; die wahrgenommene Benutzerfreundlichkeit (PEOU) der Grad, zu dem sie die Benutzung für frei von Anstrengung hält. In zwei Untersuchungen mit 152 Personen erreichen die beiden Skalen aus je sechs Aussagen Reliabilitäten von 0,98 und 0,94. Die Nützlichkeit hängt stärker mit der Nutzung zusammen als die Benutzerfreundlichkeit, etwa mit 0,85 gegenüber 0,59 für die vorhergesagte Nutzung. Wird die Nützlichkeit kontrolliert, verschwindet der Beitrag der Benutzerfreundlichkeit fast vollständig. Davis liest das vorläufig als Kette: Benutzerfreundlichkeit wirkt über die Nützlichkeit auf die Nutzung. Normwerte berichtet er nicht.

\abbildung{abb-3-1-tam.png}{0.85}{abb:tam}%
{Die von Davis (1989) berichteten Zusammenhänge und die von Shroff et al. (2011) geprüften Pfade}%
{Die von Davis (1989) berichteten Zusammenhänge und die von Shroff et al. (2011) geprüften Pfade (eigene Darstellung nach Davis, 1989, und Shroff et al., 2011)}

Shroff et al. (2011) übertragen das Modell auf ein E-Portfolio-System, das 72 Studierende ein Semester lang genutzt hatten. Abbildung~\ref{abb:tam} stellt beide Untersuchungen nebeneinander: Abgesichert sind bei Shroff et al. nur die beiden Pfade, die von der Benutzerfreundlichkeit ausgehen, darunter der zur Nützlichkeit mit 0,71 als stärkster. Der gemeinsame Nenner ist schmal, liegt aber genau an der Stelle, auf die es hier ankommt. Stichprobe und Einrichtung sind klein, die Wirkungsrichtung setzen beide Modelle voraus, statt sie zu prüfen.

Für die vorliegende Arbeit folgt daraus eine Festlegung vor jeder Evaluation: Ist die Benutzerfreundlichkeit der Hebel auf die wahrgenommene Nützlichkeit, entscheidet sie darüber, ob sich ein freiwillig gefülltes Portal füllt, und ist eine Frage des Entwurfs, nicht der Nacharbeit. Das Modell liefert zudem das Messinstrument für Kapitel~\ref{kap:test}; welche Merkmale Benutzerfreundlichkeit herstellen, lässt Davis offen.

\section{Usability und Gestaltungsprinzipien}
\label{sec:usability}

Gebrauchstauglichkeit ist nach Brooke (1996) nur in Bezug auf einen Nutzungskontext bestimmbar; die Kriterien dieses Abschnitts gelten daher für ein institutionelles Portfolio-Portal. Das Verfahren, eine Oberfläche ohne Testpersonen zu beurteilen, geht auf Nielsen und Molich (1990) zurück: Mehrere Begutachtende prüfen unabhängig gegen eine kleine Liste von Prinzipien. In vier Versuchsreihen fand eine Person im Mittel 20 bis 51 Prozent der bekannten Probleme, fünf Personen rund zwei Drittel; daraus stammt die Empfehlung von drei bis fünf Begutachtenden und die Ergänzung durch nutzerzentrierte Verfahren wie das laute Denken (Nielsen, 1994b). Die heute geführte Liste der zehn Heuristiken geht auf eine Faktorenanalyse von 249 Problemen zurück (Nielsen, 1994a, 2024).

Die Heuristiken sind Faustregeln und keine Abnahmekriterien (Nielsen, 2024). Tabelle~\ref{tab:heuristiken} leitet je Heuristik ein Kriterium für ein Portfolio-Portal ab; diese Spalte ist eine Ableitung der vorliegenden Arbeit und der Maßstab für Abschnitt~\ref{sec:bewertung-test}. Die rechte Spalte enthält bereits das Ergebnis der Referenzanalyse aus Abschnitt~\ref{sec:verwandte-arbeiten}.

\begingroup
\footnotesize
\setlength{\tabcolsep}{4pt}
\renewcommand{\arraystretch}{1.05}
\setlength{\LTcapwidth}{\textwidth}
\begin{longtable}{@{}L{0.9cm}L{3.6cm}L{6.6cm}L{3.3cm}@{}}
	\caption[Die zehn Heuristiken, das abgeleitete Kriterium und der Befund der Referenzplattform]{Die zehn Heuristiken, das abgeleitete Kriterium und der Befund der Referenzplattform (eigene Darstellung nach Nielsen, 2024; heuristische Analyse Mai 2026)}
	\label{tab:heuristiken}\\
	\toprule
	\textbf{Nr.} & \textbf{Heuristik} & \textbf{Abgeleitetes Kriterium} & \textbf{Referenzplattform} \\
	\midrule
	\endfirsthead
	\toprule
	\textbf{Nr.} & \textbf{Heuristik} & \textbf{Abgeleitetes Kriterium} & \textbf{Referenzplattform} \\
	\midrule
	\endhead
	\bottomrule
	\endfoot
	H1 & Sichtbarkeit des Systemzustands & Zustand eines Vorgangs und wirkende Einschränkung sind ablesbar & teilweise: keine Brotkrume, Auswahl nur im Bedienelement \\
	H2 & Übereinstimmung mit der realen Welt & Benennungen folgen dem Vokabular der Hochschule & erfüllt \\
	H3 & Kontrolle und Freiheit & Einschränkung einzeln und vollständig rücknehmbar & teilweise: kein Zurücksetzen \\
	H4 & Konsistenz und Standards & Gleichartiges folgt demselben Aufbau & erfüllt \\
	H5 & Fehlervermeidung & Folgenreiche Schritte werden verhindert, bevor sie erklärt werden müssen & nicht beurteilbar \\
	H6 & Wiedererkennen statt Erinnern & Ordnungsmerkmale werden sichtbar angeboten & teilweise: keine Schlagwortübersicht \\
	H7 & Flexibilität und Effizienz & Arbeit ohne Weg über alle Ordnungsebenen erreichbar & nicht erfüllt: kein Suchfeld, Filter nicht kombinierbar \\
	H8 & Ästhetisches, minimalistisches Design & Gestaltung ordnet sich den Arbeiten unter & erfüllt \\
	H9 & Hilfe beim Beheben von Fehlern & Zurückgewiesene Eingabe nennt Grund und nächsten Schritt & nicht beurteilbar \\
	H10 & Hilfe und Dokumentation & Erklärung eines Vorgangs im System auffindbar & teilweise: keine Hilfeseite \\
\end{longtable}
\endgroup

Die Auffindbarkeit, nach der die zweite Teilfrage fragt, führen die Heuristiken nicht als eigenes Thema, und auch der SUS-Bogen erfasst sie nicht (Lewis \& Sauro, 2018). Die vorliegende Arbeit setzt sie deshalb aus drei Kriterien zusammen: Die Systematik des Bestandes liegt als auswählbare Merkmale vor (H6), jede Einschränkung ist einzeln und vollständig rücknehmbar (H3), und die wirkende Einschränkung bleibt erkennbar, weil sich ein leeres Ergebnis sonst nicht von einem leeren Bestand unterscheiden lässt (H1).

Daneben treten die Web Content Accessibility Guidelines 2.2 (W3C, 2023) als Standard mit teils rechnerisch prüfbaren Schwellen. Einschlägig sind für das Portal zwölf Erfolgskriterien, darunter Textkontrast von 4,5:1 (1.4.3), 3:1 für Bedienelemente (1.4.11), keine allein farbliche Information (1.4.1), Umbruch bei 320 CSS-Pixeln (1.4.10), sichtbarer Fokus (2.4.7), Zielgröße von 24 mal 24 CSS-Pixeln (2.5.8) sowie Tastaturbedienung, Struktur und Name, Rolle und Zustand (2.1.1, 1.3.1, 4.1.2). Zielniveau ist AA.

Als drittes Instrument dient die System Usability Scale (SUS) nach Brooke (1996): zehn abwechselnd gepolte Aussagen, die einen Wert von 0 bis 100 ergeben. Brooke nennt keinen Vergleichswert; diesen liefert eine Notenskala aus 241 Untersuchungen, auf der 68 die Mitte der Note C und den Median markiert (Sauro \& Lewis, 2016, zitiert nach Lewis \& Sauro, 2018). Ein unbequemer Befund gehört an diese Stelle: Davis (1989) fand, dass Flexibilität nicht immer mit Benutzerfreundlichkeit einhergeht, was ein Werkzeug mit gestalterischer Freiheit unmittelbar betrifft.

\section{Verwandte Arbeiten und Referenzplattformen}
\label{sec:verwandte-arbeiten}

Schon für ein Inhaltsmodell des studentischen E-Portfolios sieht Mahasneh (2020) keine Vorarbeit; für hochschulische Portale gilt das nach Einschätzung der vorliegenden Arbeit erst recht. Die Literatur ordnet dagegen die Wege zu einem solchen System: Lorenzo und Ittelson (2005) unterscheiden Eigenentwicklung, quelloffene Systeme, Einkauf und gängige Werkzeuge. Die Eigenentwicklung liefert genau das Gewünschte, verlangt aber Zeit, Fachkompetenz und dauerhaft gebundenes Personal, worauf Abschnitt~\ref{sec:ausblick} zurückkommt. Freie Werkzeuge erlauben freiere Portfolios, setzen aber Kenntnisse im Erstellen von Webseiten voraus, und Shroff et al. (2011) nennen bestehende Systeme starr in der Gestaltung ihrer Inhalte. Gestalterische Freiheit ohne solche Kenntnisse ist damit die Antwort auf ein benanntes Problem.

\abbildung{abb-3-2-hfg-filter.png}{0.8}{abb:hfg}%
{Studiengangsübersicht der Referenzplattform}%
{Studiengangsübersicht der Referenzplattform: eine Auswahlliste als vollständiger Filterapparat (Bildschirmaufnahme portfolio.hfg-gmuend.de, abgerufen am 05.09.2026)}

Als Referenz dient das Portfolio-Portal der Hochschule für Gestaltung Schwäbisch Gmünd (Hochschule für Gestaltung Schwäbisch Gmünd, o.\,J.), das Arbeiten aus vier Bachelor-Studiengängen ohne Anmeldung zeigt. Gewählt wurde es auf Empfehlung des Erstprüfers, ohne systematisches Auswahlverfahren; die Befunde stützen keine Aussage über Portfolio-Portale allgemein. Die beiden Verfasser begutachteten im Mai 2026 fünf Klassen von Seiten nach Nielsen und Molich (1990) und bewerteten vierstufig, mit „nicht beurteilbar“ für alles hinter der Anmeldung. Tabelle~\ref{tab:heuristiken} fasst das Ergebnis zusammen: drei Heuristiken erfüllt, vier teilweise, zwei nicht beurteilbar, H7 nicht erfüllt.

Abbildung~\ref{abb:hfg} zeigt den Kern des Negativbefunds. Unter dem Seitentitel steht eine einzige Auswahlliste, und sie ist der gesamte Filterapparat; ein Suchfeld fehlt, und welche Einschränkung wirkt, sagt allein das Bedienelement. Semester und Modul erscheinen an jedem Projekt, sind aber keine auswählbaren Achsen. Erschlossen ist der Bestand über Studiengang, Schlagwort und Person, wobei die Person eine Sammelseite ohne Biografie bleibt. Daraus folgen die eigenen Anforderungen an Suche, verbindbare Filter und eine Profilseite (Abschnitt~\ref{sec:funktionale-anforderungen}).

Die Reichweite ist begrenzt: zwei statt drei bis fünf Begutachtende, eine stichprobenartige Begehung, keine Messung von Ladezeiten oder Kontrasten. Da die Außensicht zugleich die Sicht der Hauptzielgruppe ist, mindert die fehlende Anmeldung die Aussage zur öffentlichen Auffindbarkeit nicht; für den Einreichungsweg fehlt dagegen jede Vorlage.
